\documentclass[../main.tex]{subfiles}

\begin{document}
\section[Event detection]{Event detection}
\label{sec:event-detection}

Intuitively to detect that a rule application creates some orbit one would
compare the graph before and after its application. However comparing the
members of a rule against each other allow to statically detect topological
events. Such method presents the advantage to be independent of the object
affected by the rule.

\begin{figure}[H]
  \centering
  \includegraphics[width=\columnwidth]{triangulation}
  \caption{Rule $L \longrightarrow R$ of face triangulation}
  \label{fig:rule-triangulation}
\end{figure}


\begin{figure}[H]
  \centering
  \begin{subfigure}{.45\columnwidth}
    \centering
    \includegraphics[width=.7\columnwidth]{house-triangulated-left-match}
    \caption{}
    \label{fig:house-triangulation-left-match}
  \end{subfigure}
  \hfill
  \begin{subfigure}{.45\columnwidth}
    \centering
    \includegraphics[width=.7\columnwidth]{house-triangulated-right-match}
    \caption{}
    \label{fig:house-triangulation-right-match}
  \end{subfigure}
  % \hfill
  % \begin{subfigure}{.3\columnwidth}
  %   \centering
  %   \includegraphics[width=\columnwidth]{house-triangulated-right-match-edges}
  %   \caption{}
  %   \label{fig:house-triangulation-right-match-edges}
  % \end{subfigure}
  \caption{%
     \subref{fig:house-triangulation-left-match} - Graph $G$ - before triangulation.%
     \subref{fig:house-triangulation-right-match} - Graph $H$ - after triangulation.%
    % \subref{fig:house-triangulation-right-match-edges} - Edges created by triangulation.
   }
  \label{fig:before-after-house-triangulation}
\end{figure}

\subsection[Creation]{Creation}
\label{sec:creation}

In a rule a node matches one or more darts. So when a rule creates a node it
also creates the corresponding darts. If the right hand side of a rule contains
a created orbit, then the application of this rule on any object creates one or
more orbits. For example, in figure \ref{fig:rule-triangulation},
the face triangulation rule creates the orbit
$n_{2}$·$\orbit{\alpha 0, \alpha 2}$ in the right hand side. Therefore the
application of this rule on the square face, in figure
\ref{fig:house-triangulation-left-match}, creates four edges (emphasis in purple
in figure \ref{fig:house-triangulation-right-match}).


\begin{definition}[Orbit creation]
  \label{def:orbit-creation}
  Let $r : L \longrightarrow R$ be a rule, $\orbit{o}$ an orbit type and $n \in R$ a created
  node. An orbit $R\orbit{o}(n)$ is created if it is a complete orbit and if
  all of its nodes are created.
\end{definition}

%then and for all darts
%   $d$ instance of $n$, $H\orbit{o}(d)$ is created in $H$.
%
% \agnes{En fait il faudrait définir ce qu'est une orbite complète (si ce n'est
%   pas déjà fait dans la section précédente) et une orbite créée.}

% \agnes{La def peut pas contenir une proposition. Il faut y mettre seulement la
%   condition syntaxique sur la règle dans la def. Ensuite tu as 2 possibilités.
%   (1) faire une proposition, mais tu devras au minimum donner une idée de la
%   preuve. (2) expliquer à la suite de la def dans texte que c'est une condition
%   nécessaire et suffisante pour la règle et l'expliquer de manière plus
%   intuitive, en t'appuyant sur l'exemple.} Definition \ref{def:orbit-creation}
% gives the necessary and sufficient conditions under which we can confirm a rule
% will create some orbit.
Definition \ref{def:orbit-creation} gives the necessary and sufficient
conditions under which we can confirm a rule will create an orbit.

Given the face triangulation rule in figure \ref{fig:rule-triangulation} and its
instanciation in figure \ref{fig:before-after-house-triangulation} we will
follow the creation event of a vertex of orbit type
$\orbit{\alpha 1, \alpha 2}$. This rule creates two nodes in $R$, namely,
$n_{1}$ and $n_{2}$. Starting with $R\orbit{\alpha 1, \alpha 2}(n_{1})$ we
obtain an orbit containing nodes $n_{0}$ and $n_{1}$. As this orbit contains the
preserved node $n_{0}$ we cannot consider the set of darts it matches as fully
created. On the contrary the orbit $R\orbit{\alpha 1, \alpha 2}(n_{2})$ only
contains $n_{2}$. As this orbit only has created nodes it is fully created. Thus
we can state that the application of the face triangulation rule on the square
in figure \ref{fig:house-triangulation-left-match} will create the vertex in
blue in figure \ref{fig:house-triangulation-right-match} (i.e. the orbit
$H\orbit{\alpha 1, \alpha 2}(9)$ where $9$ is a dart matched by $n_{2}$.)

\subsection[Split]{Split}
\label{sec:split}

\maxime{intuition} As an orbit matches a set of darts an orbit split creates
disjointed sets from a former set of darts. An orbit split can occur in two
manners either explicitly or implicitly.

In the first case a set of nodes belong to a common orbit and are separated into
multiple orbits. An explicit split results in as much sets of darts as orbits.

In the second case a an orbit has a specific implicit link that is
\emph{rewritten} to a dimension not present in the orbit type of interest. An
implicit split results in two or more disjointed subsets of darts all matched
by the same orbit. Although some $\alpha$-links may appear to connect darts
subsets within a G-map they do not. The reason is the dimensions of these link
do not belong to the orbit.

\begin{definition}\label{def:split}
  Let a rule $r : L \longrightarrow R$, a G-map $G$, an orbit type $o$, a node
  $n$ where $n \in L \cap R$ and an application $G \longrightarrow^{r, m} H$.

  There are two types of orbit split:
  \begin{itemize}
    \item An orbit is explicitly splitted if two preserved nodes belonging to
          the same orbit in $L$ do not belong the same orbit in $R$, i.e. there
          exists a node $n' \in L \cap R$ such that
          $L\orbit{o}(n) = L\orbit{o}(n')$ and
          $R\orbit{o}(n) \cap R\orbit{o}(n') = \emptyset$. Then and for all $d$
          and $d'$, respectively instances of $n$ and $n'$, $d\orbit{o}(n)$ in
          $G$ is splitted into $d\orbit{o}(n)$ and $d'\orbit{o}(n)$ in $H$.

    \item An orbit is implicitly splitted if for all nodes within
          $L\orbit{o}(n)$ there exists a node $n'$ whose $k$-th dimension in
          label $l_{L}(n')$ is a dimension of $\orbit{o}$ and for all nodes $n'$
          in $R\orbit{o}(n)$, the $k$-th dimension of $l_{R}(n')$ is rewritten
          either to ``\_'' or to a dimension not in $\orbit{o}$. Then and for
          all $d$ matched by $n$ there exists $d'$, with $d \neq d'$, also
          matched by $n$ such that
          $d\orbit{o}(n) \cap d'\orbit{o}(n) = \emptyset$ in $H$.
   \end{itemize}
\end{definition}

\maxime{exemples : INSERT EXAMPLES HERE}

In order to proceed with our exemple let's start with a case of implicit split.

Given the face triangulation rule figure \ref{fig:rule-triangulation} we will
follow the implicit split event of a face of orbit type
$\orbit{\alpha 0, \alpha 1}$. The orbit $L\orbit{\alpha 0, \alpha 1}(n_{0})$
comprises node $n_{0}$ and matches a single set of darts which form a face
orbit. In this exemple nodes have two dimensions (or implicit links) in their
label. For node $n_{0}$ first implicit link is $\alpha 0$ and second implicit
link is $\alpha 1$ both in orbit type $\orbit{\alpha 0, \alpha 1}$. However the
orbit $R\orbit{\alpha 0, \alpha 1}(n_{0})$ comprises nodes $n_{0}$, $n_{1}$,
$n_{2}$ and not one of them has a second implicit link whose dimension is in
$\orbit{\alpha 0, \alpha 1}$. Under this condition we can affirm the rule figure
\ref{fig:rule-triangulation} splits a face into two or more faces. In figure
\ref{fig:house-triangulation-right-match} the orbit matches four disjointed
subsets of darts and each subset form a face orbit. The black and red dotted
cell highlights a face resulting from the split.


\maxime{remarques : INSERT REMARKS HERE}

\subsection[Merge]{Merge}
\label{sec:merge}

% \begin{figure}[H]
%   \centering
%   \includegraphics[scale=0.8]{deleteedge.pdf}
%   \caption{Rule $L \longrightarrow R$ of edge suppression}
%   \label{fig:rule-delete-edge}
% \end{figure}


\end{document}

%%% Local Variables:
%%% mode: latex
%%% TeX-master: "../main"
%%% End:
